% !PRINTED 355
\weberpage{355}

%<<CONT
is a rational integer divisible by $\mathfrak{m}$, then it must be divisible by $a$; if we set it $= ac$, then $c$ is a rational integer and
\weq{6}{D = b^2 - 4ac,}
\weq{7}{b^2 \equiv D \pmod{4a},}
and if the congruence (7) is satisfied, then every number $b'$ congruent to $b$ with respect to the modulus $2a$ satisfies the same congruence.

From this we conclude:

5. \gesperrt{Theorem. Every primary ideal $\mathfrak{m}$ coprime to $Q$ whose norm $= a$ yields us, through (4), one and only one root $b$ of the congruence (7), provided that not individual numbers, but classes of numbers with respect to the modulus $2a$, are regarded as roots.}

The two numbers
\weq{8}{a, \quad \frac{b + \sqrt{D}}{2}}
cannot, if $a$ is relatively prime to $Q$, both be divisible by one and the same rational prime; for otherwise
\[ \frac{b + \sqrt{D}}{2} - \frac{b - \sqrt{D}}{2} = \sqrt{D} \]
would have to be divisible by $p$; hence $D$ would be divisible by $p^2$, and $p$ would have to be a divisor of $Q$, contrary to the assumption.

For an odd $p$ this is evident, but for $p = 2$ it would follow:
\[ b \equiv 0 \pmod{2}, \quad b^2 - D \equiv 0 \pmod{16}, \]
and hence
\[ D \equiv 0, 4 \pmod{16}. \]
Thus $D/4$ would still be a discriminant and 2 would have to divide $Q$. With this it is proved:

6. Theorem. \gesperrt{If $a$ is a natural number coprime to $Q$ and $b$ is a root of (7), then one obtains a definite primary ideal $\mathfrak{m}$ as the greatest common divisor of the two numbers (8).}

Now we prove:

7. \gesperrt{Theorem. Every number $\mu$ of the order $[Q]$ that is divisible by the primary ideal $\mathfrak{m}$ and coprime to $Q$ is representable in the form}
\weq{9}{\mu = ax + \frac{b + \sqrt{D}}{2} y}
\gesperrt{once and only once, if $x, y$ are rational integers. If $x, y$ are regarded as vari-}
%CONT>>
